
\section{价格行为}
截至13:27实时价260.37元，较前收涨10\%，成交额约225.34亿元。2026年1月5日至7月13日复权序列显示，60日涨幅62.12\%，20日涨幅8.80\%，距年度高点回撤14.73\%。MA5/MA10/MA20/MA60约为242.95/240.23/249.45/218.62元。技术趋势强，但短期波动和估值拥挤同时上升。

\begin{exhibitbox}[表6-1\quad 二级市场结构]
\begin{tabularx}{\exhibitboxwidth}{L{3.2cm}X L{3.0cm}X}
\textbf{观察项} & \textbf{证据} & \textbf{判断} & \textbf{动作} \\
成交与换手 & 7月9日成交额277.86亿元、换手8.10\% & 高流动性、高博弈 & 不追单日加速 \\
龙虎榜 & 7月2/3/9连续上榜；7月9机构口径净卖约8.70亿元 & 机构多空分歧 & 等新信息而非追席位 \\
深股通/席位 & 深股通与机构席位同时出现 & 资金方向不单一 & 不把单一席位等同长期资金 \\
融资 & 7月13日融资余额140.43亿元，净偿还9.34亿元 & 杠杆拥挤与波动放大 & 回撤时控制风险预算 \\
相对表现 & 60日涨幅62.12\% & 已明确启动而非低位布局 & 关注分母能否追上价格 \\
\end{tabularx}
\sourcenote{sources/market-20260714/dragon\_tiger\_20260701\_20260714.json；eastmoney\_margin\_20260714.html；data/technical\_snapshot\_20260714.md}
\end{exhibitbox}

\section{机构票还是游资票}
Q1正式股东表显示，HKSCC持股3.79\%，三户保险账户合计3.05\%，睿远、华泰柏瑞沪深300ETF和摩根新兴动力合计2.03\%，实控人家族合计33.25\%。因此它有真实机构底仓，但这些持仓截至2026年3月31日，不能直接当作7月实时买入。

7月14日龙虎榜进一步说明边际资金是混合的：机构专用席位买入20.50亿元、卖出16.13亿元，净买4.37亿元；深股通买入9.31亿元、卖出12.57亿元，净卖3.26亿元；龙虎榜成交约70.79亿元，占全天成交29.67\%。结合7月9日机构净卖约8.70亿元，资金方向存在强烈轮动。

结论：东山精密是“机构底仓+主动交易强化”的高成交趋势票，不是纯游资票，也不是低位安静吸筹的纯机构票。机构专用席位不能穿透到具体基金，深股通不能等同于某一家外资，融资余额也属于杠杆资金而不是长期基金。

\begin{exhibitbox}[表6-2\quad 资金层级与可确认程度]
\begin{tabularx}{\exhibitboxwidth}{L{3.0cm}X L{3.0cm}X}
\textbf{资金层} & \textbf{可确认事实} & \textbf{性质} & \textbf{结论} \\
长期机构/基金 & Q1保险、睿远、ETF、摩根和HKSCC持股 & 季度末静态底仓 & 有机构基础，不代表7月仍未变 \\
机构专用席位 & 7月14日净买4.37亿元，2买4卖 & 主动交易资金 & 机构参与但分歧大 \\
深股通 & 7月14日净卖3.26亿元 & Stock Connect aggregate & 与内资机构方向相反 \\
活跃营业部 & 国泰海通上海、武汉紫阳东路等席位上榜 & 席位交易，非实名基金 & 游资/量化风格可能参与，但不能确认身份 \\
融资资金 & 7月13日余额140.43亿元、净偿还9.34亿元 & 杠杆资金 & 放大波动，不等于长期看多 \\
\end{tabularx}
\sourcenote{2026年一季度报告；data/capital\_structure\_20260714.json；sources/market-20260714/}
\end{exhibitbox}

\section{交易风格}
标的分类为\textbf{趋势波段}，不是\textbf{龙头战法}。它具备板块龙头叙事和极高成交，但7月的涨停、大跌、再涨停与机构对轰说明筹码交换激烈。支撑参考232、221.68和204元；阻力参考261.32和274.50元。只有当价格行为与Q2/H1业绩证据同向时，突破才具有研究价值。


\section{资金风格评分算法}
机构底仓分为81.07分，主动交易分为94.11分，综合分为88.24分。机构底仓分使用季度末可识别保险/公募/ETF与HKSCC；主动交易分使用龙虎榜参与度、上榜重复度、机构方向冲突和融资拥挤度。该评分的结论是“机构底仓+主动交易强化”，而不是对具体基金或游资身份的猜测。
